\section{Introduction}
Face à la raréfaction des ressources en eau, aux effets du changement climatique et à la dégradation des sols, notamment la salinisation
qui compromet la productivité de nombreuses terres agricoles au Sénégal, l’amélioration des pratiques agricoles constitue un défi primordial.
Dans ce contexte, l’agriculture de précision vise à optimiser la gestion des ressources grâce aux capteurs, aux données météorologiques, aux 
technologies numériques et aux outils d’aide à la décision. Elle permet d’adapter les apports en eau et en intrants aux besoins réels des cultures
plutôt que de recourir à des méthodes fixes et empiriques, contribuant ainsi à améliorer les rendements et à préserver les ressources hydriques.